\subsection{控制、反馈与算子视角}

Koopman 式材料属于本章最先进的部分。
它的业务价值，在于始终把反馈保留在视野之中。
行为系统不是开环的。
它们会观察、行动、接收误差，然后再次行动。

\begin{discussion}
基于算子的视角提供了一条可能的形式化升级路径，
因为它们尝试用变换后的线性对象来表示非线性的受控动力学。
这并不意味着人类问题因此就被简化为已解决的问题。
它真正提供的，是后续附录的一条严肃方向：
在不假装控制律不存在的前提下，如何研究受控的非线性转变系统？
\end{discussion}